\documentclass[conference]{IEEEtran}
\usepackage{amsmath,amssymb,booktabs,array,algorithm,algpseudocode,microtype,graphicx,pgfplots,placeins}
\usepackage[hidelinks]{hyperref}
\pgfplotsset{compat=1.18}
\usepgfplotslibrary{groupplots}
\title{Lightweight Trap-Based Free-Rider Screening with Fresh-Certificate Memory in Federated Learning}
\author{\IEEEauthorblockN{Aashnan Rahman}
\IEEEauthorblockA{Department of Computer Science and Engineering\\
Islamic University of Technology\\MSc Thesis}}
\begin{document}
\maketitle
\begin{center}\small
Paper v5: detector v15.4, created 15 October 2026.\\
\ifpaperresultsSingle-seed development evaluation; author review required
\else Blank results template: not an experimental finding\fi
\end{center}

\begin{abstract}
Free-rider screening must distinguish fabricated updates from heterogeneous
honest learning without relying on attacker identities or client-reported work.
We describe a server-observable challenge protocol combining same-trap robust
norm and tensor-profile measurements with fresh separation certificates.
Overlapping low-norm cohorts cannot authorize removal through concentration
alone. Statistical certificates persist across ambiguous cycles, but confirmation
requires a fresh certificate of the same component on both traps. Two consecutive
clean cycles clear memory; exact sendback has a separate consecutive-cycle path.
This revision provides reproducible templates for sixteen CIFAR-10 executions:
IID and Dirichlet non-IID partitions, FR1--FR4, 30\% and 40\% free riders, seed 42.
Numerical results and curves are withheld until every execution completes and
its logs validate. The protocol is an empirical geometric heuristic, not a
calibrated false-positive bound or proof of training. A development seed cannot
establish generalization to new partitions or adaptive adversaries.
\end{abstract}
\begin{IEEEkeywords}
federated learning, free riders, robust statistics, challenge models, non-IID data
\end{IEEEkeywords}

\section{Introduction and Related Work}
Federated learning aggregates locally computed contributions while retaining
training data at clients. Model averaging is a foundational example of this
approach~\cite{mcmahan2017communication}. Participants that consume the shared
model without performing expected training create a free-rider problem.
Constructed updates can be nonzero and history-dependent, making equality
testing alone insufficient. Prior work studies fabricated contributions and
defenses~\cite{lin2019freeriders,fraboni2021freerider}.

Heterogeneous honest data complicates attribution: persistent low update norms
can reflect legitimate learning rather than misconduct. Conversely, a fabricated
response may be anomalous repeatedly without satisfying adjacent-cycle
confirmation. We therefore distinguish nomination, fresh separation evidence,
temporary review, and permanent removal. This revision describes v15.4, not
the older v8 anchor/penalty method. Its contributions are a strict norm gate,
component-specific certificate memory, bounded review independent of memory,
and separate accounting of suspicion, removal, exclusion and utility. All
scheduled outcomes must be reported, not only successful attacks or seeds.

\section{Threat Model and Simulator Attacks}
The simulator uses 100 attributable clients and full participation among clients
not permanently removed. The detector uses received/returned states, public IDs
and its own schedule. It does not consume attacker labels, attack type, private
loss, sample count, telemetry or an expected removal count. The simulator knows
attack configuration to fabricate responses; evaluation labels are appended
after decisions. This API boundary does not imply that the whole simulation
application is unaware of ground truth.

Attackers receive no trap-status flag. Their evaluated implementations use only
their own received-model histories and local randomness. This is not a
cryptographic unpredictability claim against a seed-aware adversary. Collusion,
deliberate evasion, partial/on-off training, identity changes, dropouts, secure
aggregation and populations without an honest majority are outside the study.

The four behaviors are simulator settings, not detection signatures:
\begin{itemize}
\item \textbf{FR1: current sendback.} Return the currently received model unchanged,
not the previous-model playback described in the v8 paper.
\item \textbf{FR2: bounded random update.} Add independent uniform noise in
$[-10^{-3},10^{-3}]$ to each floating coordinate of the received model.
\item \textbf{FR3: received-history averaging.} Return the mean of up to five
received models including the current one. Traps enter the same history.
\item \textbf{FR4: received-trajectory imitation.} Average up to five successive
received-model differences and add Gaussian noise whose norm is 10\% of the
predicted norm. Cold start uses Gaussian scale $10^{-4}$; implementation norm
normalization has numerical floor $10^{-6}$.
\end{itemize}
The detector does not reconstruct these history windows. Different windows
remain a required follow-up experiment rather than an established robustness
claim.

\section{Proposed Methodology}
\subsection{Warm-Up and Challenge Coverage}
Rounds 1--10 send the ordinary global model. Sendbacks are observed, but warm-up
cannot remove clients. At each dense-cycle start, create two independent traps
around that cycle's frozen global model. For floating state tensor $k$,
\begin{equation}
 \tau_{a,k}=w_k+\gamma\max(\operatorname{sd}(w_k),\epsilon_t)\xi_{a,k},
 \quad\xi_{a,k}\sim\mathcal N(0,I),
\end{equation}
where $a\in\{A,B\}$, $\gamma=0.01$, $\epsilon_t=10^{-4}$. Traps A/B are
independent, not complementary plus/minus perturbations.

Ten shuffled groups give each client one response to each trap during rounds
11--20. Later cycles form up to ten nonempty groups from monitored and ordinary
pools separately, with proportional nearest rounding and half ties downward.
Monitored clients receive alternating traps in all ten rounds; ordinary clients
receive each trap once. A single response position per client per trap is
selected before scores are observed. Extra responses cannot become additional
confirmation strikes in the cycle. Fresh traps and grouping are used each cycle.

Only bounded-review clients automatically receive dense monitoring. A weak C
label does not impose indefinite probing or quarantine. When review/probation
are empty, an all-client single-trap audit occurs every ten rounds. It can
initiate bounded follow-up, but cannot confirm removal or clear certificate
memory.

\subsection{Same-Trap Norm and Tensor Profile}
Define $\Delta_i=w_i^{\rm ret}-w_i^{\rm send}$ and $r_i=\|\Delta_i\|_2$ over
floating tensors. Profile coordinate
$p_{i,k}=\|\Delta_{i,k}\|_2/\sum_j\|\Delta_{i,j}\|_2$ describes tensor-wise
norm allocation. A denominator below $10^{-6}$ yields a zero profile; that
convention is not the exact-sendback criterion.

Each selected response is compared with all other available active responses
to the same trap, without trusted-R or ground-truth filtering. For peer norm
median $m_i$ and MAD $d_i$,
\begin{equation}
 z_i^n=\frac{r_i-m_i}{d_i/\Phi^{-1}(0.75)}.
\end{equation}
Zero MAD produces an undefined score, not an arbitrary MAD floor. Normal
consistency is not a Gaussian-update assumption. Profile distance is Euclidean
distance from the peer coordinate-wise median profile. Its robust score uses
peer distances about that same reference. Cosine is not collected or used.

\subsection{Nomination Is Not Removal Evidence}
Existing diagnostic screening uses $D_i=z_i^p-z_i^n$, its A/B values and their
average. Each view uses $\lceil\sqrt N\rceil$ histogram bins, an Otsu
between-class-variance boundary, and 1,000 whole-client bootstrap resamples with
a fixed reproducibility seed. All three values above the upper 99\%
boundary-stability bands nominate suspicion; all below the lower bands are
unflagged, otherwise C. An unsupported majority split falls back to norm signs
about same-trap medians. These nominations cannot alone remove clients. The
99\% band is not a 1\% false-removal guarantee.

\subsection{Fresh Component Certificates}
Let $C$ be the nominated cohort and $P$ its complement for one trap. Require
$\lceil N/\lceil\sqrt N\rceil\rceil\leq |C|<|P|$. With cohort spread
$I_C=Q_{.75}(C)-Q_{.25}(C)$, the strict low-norm gate requires
\begin{align}
 Q_{.25}(P)-Q_{.75}(C)&>I_C,\\
 \min(P)-Q_{.75}(C)&>I_C.
\end{align}
An individual must also lie below the peer median. Concentration alone cannot
authorize norm removal. This separates a robust core, not necessarily every
member: the upper quartile is not the cohort maximum.

High raw profile distances retain the v15.1 gate. Let $I_P$ be peer IQR and
$g=Q_{.25}(C)-Q_{.75}(P)$. Require $g>I_C$, $g>0$, and either
\begin{equation}
 Q_{.25}(C)-\max(P)>I_C
 \quad\text{or}\quad
 \frac{I_C}{|C|-1}<\frac{I_P}{|P|-1}.
\end{equation}
The individual must exceed the peer median. The same component must certify
on both selected traps. Invalid/missing measurements cannot certify. Diagonal
scores remain diagnostic because they already contain norm evidence.

\subsection{Memory, Clearance and Exact Sendback}
Retain $M_{i,j}$ separately for norm and profile components. A fresh full A/B
certificate increments its own count once per distinct completed cycle.
Ambiguous C cycles pause without incrementing or erasing counts. Two consecutive
complete R cycles clear them. Statistical removal requires a fresh current
certificate with $M_{i,j}\geq2$. Components never combine and duplicate/older
cycles are rejected. There is no round-30/60 switch, removal quota or arbitrary
certificate age limit. Long retention nevertheless risks accumulating
correlated false evidence.

Exact current zero delta ($r_i=0$), or a fingerprint match to the immediately
previous model sent to that client, is a separate sendback observation. Sendback
on both selected traps in consecutive complete cycles can remove. Near-zero
norms alone are not exact equality, and warm-up cannot bypass this requirement.

\begin{algorithm}[t]
\caption{Component-specific fresh-certificate confirmation}
\begin{algorithmic}[1]
\Require complete A/B measurements and a distinct cycle
\State compute unchanged v15.3 nominations and strict certificates
\ForAll{active clients $i$}
 \If{fresh statistical certificate for component $j$}
  \State $M_{i,j}\gets M_{i,j}+1$; interrupt clean streak
  \If{$M_{i,j}\geq2$} \State remove $i$ \EndIf
 \ElsIf{complete R cycle}
  \State advance consecutive clean-cycle streak
  \If{two clean cycles} \State clear all $M_{i,j}$ \EndIf
 \Else \State retain counts; interrupt clean streak
 \EndIf
 \State apply separate consecutive-cycle exact-sendback rule
\EndFor
\State release unconfirmed quarantine when bounded review expires
\end{algorithmic}
\end{algorithm}

\subsection{Aggregation and Recovery}
Trap responses never aggregate. Certified S starts a two-cycle bounded review;
fresh strong evidence can renew it, but stored memory and weak C evidence cannot.
Expired unconfirmed quarantine is released. Two clean cycles rehabilitate,
followed by one extra clean probation cycle. Weak candidates can aggregate on
ordinary rounds. New quarantine/removal also excludes the decision-round update.
Eligible deltas are averaged equally, not weighted by private sample counts.
An empty eligible set leaves the global model unchanged.

Independent streams seed partitioning, attacker selection, initialization,
training, grouping, traps and attack randomness. Checkpoints retain RNG/loaders,
attacker histories, selected probes, review state and memory. Resume rejects
changed protocol/source hashes. Recovery does not create detector evidence.

\section{Experimental Design}
The intended grid is CIFAR-10, IID/non-IID, FR1--FR4, 30\%/40\% free riders:
sixteen 100-round executions with master seed 42. IID uses shuffled equal-size
partitions. Non-IID uses class-wise Dirichlet allocation with $\alpha=0.5$,
retrying to obtain nonempty client datasets. Honest clients use three local SGD
epochs, batch size 32, learning rate 0.01 and momentum 0.9. Images use fixed
channel normalization without augmentation. The CNN has three convolution/
BatchNorm blocks (32, 64, 128 channels), pooling, a 256-unit FC layer and ten logits.

Evaluation reports TP/FP/FN for permanent removal, precision, recall, F1 and
honest false-removal rate. Precision is undefined when nobody is removed.
Complete-detection round $t_{\rm all}$ exists only if all free riders are removed;
otherwise report NR, not the last observed removal. Utility, final honest
candidates, unique honest clients quarantined, quarantine client-rounds and
whole-process resources are separate. Reused identities are client-cases, not
independent statistical samples. A single seed does not support sample-SD or
cross-seed confidence intervals.

The reporting pipeline requires all sixteen completion markers, 100 rounds per
run, unique labelled removals, fixed evaluation truth, complete initial A/B
coverage, consistent active-client logs and confusion/exclusion summaries, and
matching frozen detector provenance. It neither resumes training nor selects
favorable runs. Incomplete data stay blank; completed detection failures remain.
Earlier MNIST/v8/v15.3 outcomes are not substituted.

\section{Results Templates}
\ifpaperresults
Across sixteen completed seed-42 executions, 560/560 free-rider cases and 0/1040 honest-client cases were permanently removed. Pooled removal recall is 100.00\%, precision is 100.00\%, and F1 is 100.00\%. There are 10 honest-client cases ever quarantined and 182 final honest-candidate cases. These are repeated simulation cases, not independent clients or held-out seeds. See the per-condition tables; incomplete removal remains a failure rather than a missing result.

\begin{table*}[t]
\centering\footnotesize\setlength{\tabcolsep}{4pt}
\caption{Removal performance. NR means complete removal was not achieved by round 100; -- denotes pending/undefined evidence.}\label{tab:detection}
\begin{tabular}{rllrrrrrrr}
\toprule
FR (\%) & Split & Attack & TP & FP & FN & Recall (\%) & Precision (\%) & F1 (\%) & $t_{\rm all}$ \\
\midrule
30 & IID & FR1 & 30 & 0 & 0 & 100.00 & 100.00 & 100.00 & 30 \\
30 & IID & FR2 & 30 & 0 & 0 & 100.00 & 100.00 & 100.00 & 30 \\
30 & IID & FR3 & 30 & 0 & 0 & 100.00 & 100.00 & 100.00 & 30 \\
30 & IID & FR4 & 30 & 0 & 0 & 100.00 & 100.00 & 100.00 & 30 \\
30 & Non-IID & FR1 & 30 & 0 & 0 & 100.00 & 100.00 & 100.00 & 30 \\
30 & Non-IID & FR2 & 30 & 0 & 0 & 100.00 & 100.00 & 100.00 & 30 \\
30 & Non-IID & FR3 & 30 & 0 & 0 & 100.00 & 100.00 & 100.00 & 80 \\
30 & Non-IID & FR4 & 30 & 0 & 0 & 100.00 & 100.00 & 100.00 & 30 \\
40 & IID & FR1 & 40 & 0 & 0 & 100.00 & 100.00 & 100.00 & 30 \\
40 & IID & FR2 & 40 & 0 & 0 & 100.00 & 100.00 & 100.00 & 30 \\
40 & IID & FR3 & 40 & 0 & 0 & 100.00 & 100.00 & 100.00 & 30 \\
40 & IID & FR4 & 40 & 0 & 0 & 100.00 & 100.00 & 100.00 & 30 \\
40 & Non-IID & FR1 & 40 & 0 & 0 & 100.00 & 100.00 & 100.00 & 30 \\
40 & Non-IID & FR2 & 40 & 0 & 0 & 100.00 & 100.00 & 100.00 & 30 \\
40 & Non-IID & FR3 & 40 & 0 & 0 & 100.00 & 100.00 & 100.00 & 60 \\
40 & Non-IID & FR4 & 40 & 0 & 0 & 100.00 & 100.00 & 100.00 & 30 \\
\bottomrule
\end{tabular}
\end{table*}

\begin{table*}[t]
\centering\footnotesize\setlength{\tabcolsep}{4pt}
\caption{Utility and honest-client safety. $C_H$ is final honest-candidate count, $Q_H$ unique honest clients ever quarantined, $E_H$ honest quarantine client-rounds, and $Q_T$ all clients quarantined at completion.}\label{tab:safety}
\begin{tabular}{rllrrrrr}
\toprule
FR (\%) & Split & Attack & Accuracy (\%) & $C_H$ & $Q_H$ & $E_H$ & $Q_T$ \\
\midrule
30 & IID & FR1 & 78.07 & 2 & 0 & 0 & 0 \\
30 & IID & FR2 & 77.93 & 0 & 0 & 0 & 0 \\
30 & IID & FR3 & 78.17 & 2 & 0 & 0 & 0 \\
30 & IID & FR4 & 78.03 & 7 & 0 & 0 & 0 \\
30 & Non-IID & FR1 & 75.80 & 8 & 0 & 0 & 0 \\
30 & Non-IID & FR2 & 75.58 & 2 & 0 & 0 & 0 \\
30 & Non-IID & FR3 & 75.36 & 70 & 1 & 20 & 0 \\
30 & Non-IID & FR4 & 75.78 & 7 & 0 & 0 & 0 \\
40 & IID & FR1 & 77.46 & 0 & 0 & 0 & 0 \\
40 & IID & FR2 & 77.65 & 0 & 0 & 0 & 0 \\
40 & IID & FR3 & 77.27 & 4 & 0 & 0 & 0 \\
40 & IID & FR4 & 77.36 & 8 & 0 & 0 & 0 \\
40 & Non-IID & FR1 & 75.28 & 9 & 0 & 0 & 0 \\
40 & Non-IID & FR2 & 75.43 & 4 & 0 & 0 & 0 \\
40 & Non-IID & FR3 & 74.89 & 52 & 9 & 180 & 0 \\
40 & Non-IID & FR4 & 75.14 & 7 & 0 & 0 & 0 \\
\bottomrule
\end{tabular}
\end{table*}

\begin{table*}[t]
\centering\footnotesize\setlength{\tabcolsep}{4pt}
\caption{Whole-simulator resources, not isolated detector overhead. Mean logged round time excludes subsequent checkpoint/log publication; wall time includes initialization and publication. RSS includes training, aggregation and simulated attacker histories.}\label{tab:resources}
\begin{tabular}{rllrrr}
\toprule
FR (\%) & Split & Attack & Mean round (s) & Peak RSS (MiB) & Wall time (min) \\
\midrule
30 & IID & FR1 & 88.18 & 987.07 & 152.4 \\
30 & IID & FR2 & 97.78 & 780.39 & 165.6 \\
30 & IID & FR3 & 99.87 & 1796.55 & 179.0 \\
30 & IID & FR4 & 119.07 & 1697.04 & 217.4 \\
30 & Non-IID & FR1 & 95.03 & 983.02 & 162.9 \\
30 & Non-IID & FR2 & 89.56 & 752.98 & 151.7 \\
30 & Non-IID & FR3 & 108.38 & 1807.34 & 193.5 \\
30 & Non-IID & FR4 & 95.61 & 1791.95 & 173.6 \\
40 & IID & FR1 & 80.97 & 1013.29 & 139.7 \\
40 & IID & FR2 & 77.38 & 734.54 & 131.3 \\
40 & IID & FR3 & 104.21 & 1936.14 & 198.0 \\
40 & IID & FR4 & 100.07 & 1977.06 & 193.8 \\
40 & Non-IID & FR1 & 78.42 & 1040.16 & 135.9 \\
40 & Non-IID & FR2 & 75.80 & 794.95 & 128.9 \\
40 & Non-IID & FR3 & 84.85 & 2180.93 & 160.1 \\
40 & Non-IID & FR4 & 78.79 & 2118.29 & 149.5 \\
\bottomrule
\end{tabular}
\end{table*}

\begin{figure*}[t]\centering
\begin{tikzpicture}
\begin{groupplot}[group style={group size=2 by 2,horizontal sep=1.15cm,vertical sep=1.35cm},
width=0.44\textwidth,height=3.5cm,xmin=0,xmax=100,ymin=0,ymax=1,
xlabel={Round},ylabel={Cumulative removal recall},grid=major,legend style={font=\scriptsize},
tick label style={font=\scriptsize},title style={font=\small}]
\nextgroupplot[title={30\% IID}]
\addplot[blue,solid,thick] table[x=round,y=recall] {generated/iid_FR1_p30.dat};
\addlegendentry{FR1}
\addplot[red,dashed,thick] table[x=round,y=recall] {generated/iid_FR2_p30.dat};
\addlegendentry{FR2}
\addplot[teal,dotted,thick] table[x=round,y=recall] {generated/iid_FR3_p30.dat};
\addlegendentry{FR3}
\addplot[orange,dashdotted,thick] table[x=round,y=recall] {generated/iid_FR4_p30.dat};
\addlegendentry{FR4}
\nextgroupplot[title={30\% Non-IID}]
\addplot[blue,solid,thick] table[x=round,y=recall] {generated/noniid_FR1_p30.dat};
\addlegendentry{FR1}
\addplot[red,dashed,thick] table[x=round,y=recall] {generated/noniid_FR2_p30.dat};
\addlegendentry{FR2}
\addplot[teal,dotted,thick] table[x=round,y=recall] {generated/noniid_FR3_p30.dat};
\addlegendentry{FR3}
\addplot[orange,dashdotted,thick] table[x=round,y=recall] {generated/noniid_FR4_p30.dat};
\addlegendentry{FR4}
\nextgroupplot[title={40\% IID}]
\addplot[blue,solid,thick] table[x=round,y=recall] {generated/iid_FR1_p40.dat};
\addlegendentry{FR1}
\addplot[red,dashed,thick] table[x=round,y=recall] {generated/iid_FR2_p40.dat};
\addlegendentry{FR2}
\addplot[teal,dotted,thick] table[x=round,y=recall] {generated/iid_FR3_p40.dat};
\addlegendentry{FR3}
\addplot[orange,dashdotted,thick] table[x=round,y=recall] {generated/iid_FR4_p40.dat};
\addlegendentry{FR4}
\nextgroupplot[title={40\% Non-IID}]
\addplot[blue,solid,thick] table[x=round,y=recall] {generated/noniid_FR1_p40.dat};
\addlegendentry{FR1}
\addplot[red,dashed,thick] table[x=round,y=recall] {generated/noniid_FR2_p40.dat};
\addlegendentry{FR2}
\addplot[teal,dotted,thick] table[x=round,y=recall] {generated/noniid_FR3_p40.dat};
\addlegendentry{FR3}
\addplot[orange,dashdotted,thick] table[x=round,y=recall] {generated/noniid_FR4_p40.dat};
\addlegendentry{FR4}
\end{groupplot}\end{tikzpicture}
\caption{Cumulative removal recall, seed 42. All attacks are retained, including misses; coincident curves may overlap.}
\end{figure*}

\begin{figure*}[t]\centering
\begin{tikzpicture}
\begin{groupplot}[group style={group size=2 by 2,horizontal sep=1.15cm,vertical sep=1.35cm},
width=0.44\textwidth,height=3.5cm,xmin=0,xmax=100,ymin=0,ymax=1,
xlabel={Round},ylabel={Honest false-removal rate},grid=major,legend style={font=\scriptsize},
tick label style={font=\scriptsize},title style={font=\small}]
\nextgroupplot[title={30\% IID}]
\addplot[blue,solid,thick] table[x=round,y=fpr] {generated/iid_FR1_p30.dat};
\addlegendentry{FR1}
\addplot[red,dashed,thick] table[x=round,y=fpr] {generated/iid_FR2_p30.dat};
\addlegendentry{FR2}
\addplot[teal,dotted,thick] table[x=round,y=fpr] {generated/iid_FR3_p30.dat};
\addlegendentry{FR3}
\addplot[orange,dashdotted,thick] table[x=round,y=fpr] {generated/iid_FR4_p30.dat};
\addlegendentry{FR4}
\nextgroupplot[title={30\% Non-IID}]
\addplot[blue,solid,thick] table[x=round,y=fpr] {generated/noniid_FR1_p30.dat};
\addlegendentry{FR1}
\addplot[red,dashed,thick] table[x=round,y=fpr] {generated/noniid_FR2_p30.dat};
\addlegendentry{FR2}
\addplot[teal,dotted,thick] table[x=round,y=fpr] {generated/noniid_FR3_p30.dat};
\addlegendentry{FR3}
\addplot[orange,dashdotted,thick] table[x=round,y=fpr] {generated/noniid_FR4_p30.dat};
\addlegendentry{FR4}
\nextgroupplot[title={40\% IID}]
\addplot[blue,solid,thick] table[x=round,y=fpr] {generated/iid_FR1_p40.dat};
\addlegendentry{FR1}
\addplot[red,dashed,thick] table[x=round,y=fpr] {generated/iid_FR2_p40.dat};
\addlegendentry{FR2}
\addplot[teal,dotted,thick] table[x=round,y=fpr] {generated/iid_FR3_p40.dat};
\addlegendentry{FR3}
\addplot[orange,dashdotted,thick] table[x=round,y=fpr] {generated/iid_FR4_p40.dat};
\addlegendentry{FR4}
\nextgroupplot[title={40\% Non-IID}]
\addplot[blue,solid,thick] table[x=round,y=fpr] {generated/noniid_FR1_p40.dat};
\addlegendentry{FR1}
\addplot[red,dashed,thick] table[x=round,y=fpr] {generated/noniid_FR2_p40.dat};
\addlegendentry{FR2}
\addplot[teal,dotted,thick] table[x=round,y=fpr] {generated/noniid_FR3_p40.dat};
\addlegendentry{FR3}
\addplot[orange,dashdotted,thick] table[x=round,y=fpr] {generated/noniid_FR4_p40.dat};
\addlegendentry{FR4}
\end{groupplot}\end{tikzpicture}
\caption{Honest false-removal rate, seed 42. All attacks are retained, including misses; coincident curves may overlap.}
\end{figure*}

\begin{figure*}[t]\centering
\begin{tikzpicture}
\begin{groupplot}[group style={group size=2 by 2,horizontal sep=1.15cm,vertical sep=1.35cm},
width=0.44\textwidth,height=3.5cm,xmin=0,xmax=100,ymin=0,ymax=1,
xlabel={Round},ylabel={Global test accuracy},grid=major,legend style={font=\scriptsize},
tick label style={font=\scriptsize},title style={font=\small}]
\nextgroupplot[title={30\% IID}]
\addplot[blue,solid,thick] table[x=round,y=accuracy] {generated/iid_FR1_p30.dat};
\addlegendentry{FR1}
\addplot[red,dashed,thick] table[x=round,y=accuracy] {generated/iid_FR2_p30.dat};
\addlegendentry{FR2}
\addplot[teal,dotted,thick] table[x=round,y=accuracy] {generated/iid_FR3_p30.dat};
\addlegendentry{FR3}
\addplot[orange,dashdotted,thick] table[x=round,y=accuracy] {generated/iid_FR4_p30.dat};
\addlegendentry{FR4}
\nextgroupplot[title={30\% Non-IID}]
\addplot[blue,solid,thick] table[x=round,y=accuracy] {generated/noniid_FR1_p30.dat};
\addlegendentry{FR1}
\addplot[red,dashed,thick] table[x=round,y=accuracy] {generated/noniid_FR2_p30.dat};
\addlegendentry{FR2}
\addplot[teal,dotted,thick] table[x=round,y=accuracy] {generated/noniid_FR3_p30.dat};
\addlegendentry{FR3}
\addplot[orange,dashdotted,thick] table[x=round,y=accuracy] {generated/noniid_FR4_p30.dat};
\addlegendentry{FR4}
\nextgroupplot[title={40\% IID}]
\addplot[blue,solid,thick] table[x=round,y=accuracy] {generated/iid_FR1_p40.dat};
\addlegendentry{FR1}
\addplot[red,dashed,thick] table[x=round,y=accuracy] {generated/iid_FR2_p40.dat};
\addlegendentry{FR2}
\addplot[teal,dotted,thick] table[x=round,y=accuracy] {generated/iid_FR3_p40.dat};
\addlegendentry{FR3}
\addplot[orange,dashdotted,thick] table[x=round,y=accuracy] {generated/iid_FR4_p40.dat};
\addlegendentry{FR4}
\nextgroupplot[title={40\% Non-IID}]
\addplot[blue,solid,thick] table[x=round,y=accuracy] {generated/noniid_FR1_p40.dat};
\addlegendentry{FR1}
\addplot[red,dashed,thick] table[x=round,y=accuracy] {generated/noniid_FR2_p40.dat};
\addlegendentry{FR2}
\addplot[teal,dotted,thick] table[x=round,y=accuracy] {generated/noniid_FR3_p40.dat};
\addlegendentry{FR3}
\addplot[orange,dashdotted,thick] table[x=round,y=accuracy] {generated/noniid_FR4_p40.dat};
\addlegendentry{FR4}
\end{groupplot}\end{tikzpicture}
\caption{Global test accuracy, seed 42. All attacks are retained, including misses; coincident curves may overlap.}
\end{figure*}

\else
\textit{Results pending. No quantitative findings are asserted in this permanent
blank template.}
\begin{table*}[t]
\centering\footnotesize\setlength{\tabcolsep}{4pt}
\caption{Removal performance. NR means complete removal was not achieved by round 100; -- denotes pending/undefined evidence.}\label{tab:detection}
\begin{tabular}{rllrrrrrrr}
\toprule
FR (\%) & Split & Attack & TP & FP & FN & Recall (\%) & Precision (\%) & F1 (\%) & $t_{\rm all}$ \\
\midrule
30 & IID & FR1 & -- & -- & -- & -- & -- & -- & -- \\
30 & IID & FR2 & -- & -- & -- & -- & -- & -- & -- \\
30 & IID & FR3 & -- & -- & -- & -- & -- & -- & -- \\
30 & IID & FR4 & -- & -- & -- & -- & -- & -- & -- \\
30 & Non-IID & FR1 & -- & -- & -- & -- & -- & -- & -- \\
30 & Non-IID & FR2 & -- & -- & -- & -- & -- & -- & -- \\
30 & Non-IID & FR3 & -- & -- & -- & -- & -- & -- & -- \\
30 & Non-IID & FR4 & -- & -- & -- & -- & -- & -- & -- \\
40 & IID & FR1 & -- & -- & -- & -- & -- & -- & -- \\
40 & IID & FR2 & -- & -- & -- & -- & -- & -- & -- \\
40 & IID & FR3 & -- & -- & -- & -- & -- & -- & -- \\
40 & IID & FR4 & -- & -- & -- & -- & -- & -- & -- \\
40 & Non-IID & FR1 & -- & -- & -- & -- & -- & -- & -- \\
40 & Non-IID & FR2 & -- & -- & -- & -- & -- & -- & -- \\
40 & Non-IID & FR3 & -- & -- & -- & -- & -- & -- & -- \\
40 & Non-IID & FR4 & -- & -- & -- & -- & -- & -- & -- \\
\bottomrule
\end{tabular}
\end{table*}

\begin{table*}[t]
\centering\footnotesize\setlength{\tabcolsep}{4pt}
\caption{Utility and honest-client safety. $C_H$ is final honest-candidate count, $Q_H$ unique honest clients ever quarantined, $E_H$ honest quarantine client-rounds, and $Q_T$ all clients quarantined at completion.}\label{tab:safety}
\begin{tabular}{rllrrrrr}
\toprule
FR (\%) & Split & Attack & Accuracy (\%) & $C_H$ & $Q_H$ & $E_H$ & $Q_T$ \\
\midrule
30 & IID & FR1 & -- & -- & -- & -- & -- \\
30 & IID & FR2 & -- & -- & -- & -- & -- \\
30 & IID & FR3 & -- & -- & -- & -- & -- \\
30 & IID & FR4 & -- & -- & -- & -- & -- \\
30 & Non-IID & FR1 & -- & -- & -- & -- & -- \\
30 & Non-IID & FR2 & -- & -- & -- & -- & -- \\
30 & Non-IID & FR3 & -- & -- & -- & -- & -- \\
30 & Non-IID & FR4 & -- & -- & -- & -- & -- \\
40 & IID & FR1 & -- & -- & -- & -- & -- \\
40 & IID & FR2 & -- & -- & -- & -- & -- \\
40 & IID & FR3 & -- & -- & -- & -- & -- \\
40 & IID & FR4 & -- & -- & -- & -- & -- \\
40 & Non-IID & FR1 & -- & -- & -- & -- & -- \\
40 & Non-IID & FR2 & -- & -- & -- & -- & -- \\
40 & Non-IID & FR3 & -- & -- & -- & -- & -- \\
40 & Non-IID & FR4 & -- & -- & -- & -- & -- \\
\bottomrule
\end{tabular}
\end{table*}

\begin{table*}[t]
\centering\footnotesize\setlength{\tabcolsep}{4pt}
\caption{Whole-simulator resources, not isolated detector overhead. Mean logged round time excludes subsequent checkpoint/log publication; wall time includes initialization and publication. RSS includes training, aggregation and simulated attacker histories.}\label{tab:resources}
\begin{tabular}{rllrrr}
\toprule
FR (\%) & Split & Attack & Mean round (s) & Peak RSS (MiB) & Wall time (min) \\
\midrule
30 & IID & FR1 & -- & -- & -- \\
30 & IID & FR2 & -- & -- & -- \\
30 & IID & FR3 & -- & -- & -- \\
30 & IID & FR4 & -- & -- & -- \\
30 & Non-IID & FR1 & -- & -- & -- \\
30 & Non-IID & FR2 & -- & -- & -- \\
30 & Non-IID & FR3 & -- & -- & -- \\
30 & Non-IID & FR4 & -- & -- & -- \\
40 & IID & FR1 & -- & -- & -- \\
40 & IID & FR2 & -- & -- & -- \\
40 & IID & FR3 & -- & -- & -- \\
40 & IID & FR4 & -- & -- & -- \\
40 & Non-IID & FR1 & -- & -- & -- \\
40 & Non-IID & FR2 & -- & -- & -- \\
40 & Non-IID & FR3 & -- & -- & -- \\
40 & Non-IID & FR4 & -- & -- & -- \\
\bottomrule
\end{tabular}
\end{table*}


\begin{figure*}[t]\centering
\fbox{\parbox[c][5.7cm][c]{0.94\textwidth}{\centering Cumulative removal recall\\[8pt]Blank template: populate only after all sixteen runs validate.}}
\caption{Cumulative removal recall. Panels: 30\%/40\% free riders and IID/non-IID.}
\end{figure*}
\begin{figure*}[t]\centering
\fbox{\parbox[c][5.7cm][c]{0.94\textwidth}{\centering Honest false-removal rate\\[8pt]Blank template: populate only after all sixteen runs validate.}}
\caption{Honest false-removal rate. Panels: 30\%/40\% free riders and IID/non-IID.}
\end{figure*}
\begin{figure*}[t]\centering
\fbox{\parbox[c][5.7cm][c]{0.94\textwidth}{\centering Global test accuracy\\[8pt]Blank template: populate only after all sixteen runs validate.}}
\caption{Global test accuracy. Panels: 30\%/40\% free riders and IID/non-IID.}
\end{figure*}

\fi
\FloatBarrier

\section{Cost and Limitations}
Detector features retain $O(NL)$ norm/profile summaries and a few counters per
client, plus two trap models. This is not total simulator memory: training,
aggregation and simulated FR3/FR4 histories retain additional state. Extraction
reads model coordinates; existing leave-one-out peer comparisons are roughly
quadratic in client count, and bootstrap scratch space is $O(BN)$ for $B=1000$.

Whole-process RSS and round duration cannot isolate detector overhead. A
dedicated detector-only benchmark, communication accounting and matched
no-detector controls remain required. ``Lightweight'' denotes compact
statistics without auxiliary detector training, not a deployment-scale bound.

The policy was developed using earlier seed-42 failures. Changing prevalence
to 30\% is not held-out-seed validation. Neither zero observed honest removals
nor a stable bootstrap split guarantees future error rates. Honest clients can
pass certificates, including the retained profile-density gate. Certificates
share client data and trajectories and are not statistically independent.
Minimum cohort support can miss isolated free riders; long-lived evidence can
eventually remove honest clients unless clearance occurs.

Required follow-up includes new seeds, honest-only controls, longer horizons,
different history/noise parameters, mixed/adaptive attacks, partial participation
and other architectures. Matched ablations must distinguish nomination,
certificate and memory effects. False candidate flags and temporary exclusions
must remain visible even when permanent-removal precision is high.

\section{Conclusion}
This revision specifies a strict attack-label-blind gate and component-specific
memory for intermittent fresh evidence. It supplies blank reproducible templates
rather than predicting a desired detection count. Populated findings describe
only the completed grid and require review alongside temporary honest exclusions
and single-seed limitations.
\bibliographystyle{IEEEtran}
\bibliography{refs}
\end{document}
